Perch offers three architectural advantages, each with a caveat, and one
recovery stops every earlier grant.

\paragraph{A limit over time inside the decision.} A Perch rule's window cap is
charged in the decision that checks the call. The cap theorem (Theorem~6 of
\perchdoc{verification}) states that from empty records, in any run, the accepted
charges $x_k$ at ledgers $t_k$ to a rule capped at $\lambda$ per $P$ ledgers stay
within the limit in every window:
\[
  \forall L.\ \sum_{k:\ \max(0,\,L-P) < t_k \le L} x_k \;\le\; \lambda .
\]
In oz-policy-builder the interpreter bounds each call, and the window comes from
\code{spending\_limit}, an opt-in policy on \code{transfer} only; with it removed,
the recorded dual run moved \fact{sim-oz-percall-moved-xlm}~XLM in one window
while Perch held to its cap (\S\ref{sec:same-task}). Both enforce the window
when configured with one: the difference here is where the cap is expressed
and checked, not whether a window can be enforced. \emph{Caveat:} the theorem
is about the evaluator's Lean translation; that the shell selects rules and
persists records as it assumes is argued and tested, not proved (obligations O2
and O6). On chain a record holds at most \fact{spends-per-window} spends per
window, and replacing a capped rule starts an empty window.

\paragraph{One decision point.} All the state that decides a grant's call lives
in the account: rules, grants, spending records and the grant generation
(Perch \srcpath{lib.rs:597-636}). The decision comes before authentication, so a
call refused by the evaluator invokes no verifier and writes nothing. A bad
signature instead fails during authentication. One \code{apply\_grant} installs a
grant whose rules span several contracts, and a grant rule may name a passkey or
Ed25519 key as an external signer, which the interpreter refuses in the rules it
polices. \emph{Caveat:} one point decides only what the evaluator can express
(\S\ref{sec:composition}).

\paragraph{A proved core.} The evaluator is about \fact{rule-eval-lines} lines of
\code{no\_std} Rust in a subset that Charon~\cite{charon} and Aeneas~\cite{aeneas} translate
into Lean. Kernel-checked theorems state that the translation returns what a
readable model returns on every input, without panicking or diverging
(Theorem~1); that when the decision function accepts,
each call is allowed by some rule it was given (Theorem~2); that attenuation never widens the calls a grant allows (Theorem~4); the window bound
above (Theorem~6, for fixed rules and empty initial records); and that no grant
produced by scoped lowering, nor one derived from it, allows a call to the
account or the recovery contracts excluded at install (Theorem~8). The shell
rechecks the current excluded set when recovery configuration changes. We
found no formal model in oz-policy-builder's tree (a search, not a proof of
absence); its evidence is testing, partly on the shipped Wasm.
\emph{Caveat:} the theorems are about the translation, not the compiled Wasm; the
shell is tested, not proved; attenuation cannot yet be installed; and nothing in
Perch has been audited.

\paragraph{Stopping every earlier grant.} A completed recovery advances the grant
generation, after which no earlier grant, nor one derived from it, accepts a
nonempty authorization (Theorem~10). \emph{Caveat:} this is a theorem about an
account model with no Rust counterpart; the shell's part is tested (O3).
oz-policy-builder has no account-key recovery, a difference in scope, not a
defect.
